% ============================================================
%  rolegrammar_body.tex — Main body of Paper 4
% ============================================================

% ============================================================
\section{Introduction}
\label{sec:introduction}

The \vms\ (Beinecke MS~408) has been the subject of three companion
studies by the present author. The first~\citep{jimenez2026negative}
tested nineteen substitution-based decipherment models against five
Semitic target languages and established a structured negative result:
all models converge on the same structural ceiling without producing
readable text. The second~\citep{jimenez2026twotier} demonstrated that
the manuscript's positional glyph structure is consistent with a
designed two-tier categorial code rather than a natural language in
disguise. The third~\citep{jimenez2026venusring} identified the medieval
astrological-talismanic tradition codified by \textit{Picatrix Latinus}
as the historical framework within which the manuscript's encoding
practice is situated, and established that the manuscript's global
encoding constitutes a \textit{cifra ineffabilis} in the sense of
\citet{martinezgavilan2025}.

Those three studies answered two foundational questions: \emph{what kind
of system} is the Voynich (a categorial code, not a substitution cipher)
and \emph{which tradition} does it belong to (the Picatrix-derived
talismanic pasigraphy tradition). The present paper addresses a third
question that follows naturally from the first two:

\begin{quote}
\emph{Can the combinatorial rules of the code be recovered from the
text alone, without access to the lost key book?}
\end{quote}

We show that the answer is partially affirmative. The analysis proceeds
in three stages that mirror the logical dependencies of the argument:

\begin{enumerate}[nosep]
  \item \textbf{Form encodes function} (\S\ref{sec:family_position}):
    visual glyph families correlate with word-internal positions,
    establishing that the graphic shape of a glyph is predictive of its
    combinatorial role.
  \item \textbf{Role grammar} (\S\ref{sec:role_grammar}): the
    form--function correlation resolves into a six-role combinatorial
    grammar with directional transition rules, universal across all
    manuscript sections but with section-specific glyph proportions.
  \item \textbf{Structural dictionary} (\S\ref{sec:structural_dict}):
    the grammar induces a five-slot schema for the Voynich word, each
    slot with a closed inventory and measurable section specificity,
    representing the syntactic skeleton of the lost key book.
\end{enumerate}

The progression is deliberately cumulative: each stage depends on the
results of the previous one, and none can be skipped without undermining
the next. Section~\ref{sec:discussion} discusses what the reconstructed
schema tells us---and what it does not---about the manuscript's
content. Section~\ref{sec:limitations} addresses limitations and
future work.


% ============================================================
\section{Background and prior results}
\label{sec:background}

We summarise only the results from the three companion papers that are
directly needed as premises for the present analysis. The reader is
referred to the original studies for full methodological detail.

\subsection{Corpus}

All analyses in this paper use the \sta\ 2.0 transliteration of the
\vms\ compiled by \citet{zandbergen_sta1}, consisting of 37,087 words,
157,254 glyph tokens, and 166 distinct glyph types distributed across
224 folios. Words are segmented into positional classes using the
\silpart\ pipeline introduced in the second companion
paper~\citep{jimenez2026twotier}: each glyph token is assigned to one
of six word-internal positions ($p_1$ through $p_6$) based on its
frequency profile across those positions.

\subsection{Key prior findings}

\paragraph{Two-tier positional structure.}
The second companion paper established that Voynich glyphs are not
randomly distributed across word positions but exhibit a highly
structured positional profile. Mutual information between glyph identity
and word position ($\mathrm{MI}_{\text{pos}} = 0.670$~bits for
\sta\ 2.0) is approximately double that of any natural language at
comparable inventory size. The \silpart\ generative model---which
assigns each glyph a probability distribution over six positional
slots---achieves the best generative fit among all tested models
($d = 0.445$).

\paragraph{Categorial encoding.}
The cumulative evidence from the first two papers is that the \vms\
uses a positional categorial code: each word is a sequence of glyphs
drawn from position-specific inventories, and the combinatorial rules
governing which glyph may follow which are determined by the positional
slot rather than by phonological constraints. This is the encoding type
that \citet{martinezgavilan2025} recovers under the term \textit{cifra
ineffabilis}, as discussed in the third companion
paper~\citep{jimenez2026venusring}.

\paragraph{Historical tradition.}
The third companion paper identified the medieval Picatrix tradition as
the historical context within which the Voynich's encoding practice
is situated. The present paper does not depend on the historical
identification; its results are purely structural and hold regardless
of which tradition the manuscript belongs to.


% ============================================================
\section{Form encodes function: family--position correlation}
\label{sec:family_position}

\subsection{Hypothesis}

The \sta\ 2.0 classification assigns each glyph a two-character code
whose first character denotes a \emph{visual family}: glyphs sharing
a common graphic root (e.g., all A-family glyphs share an arch-like
base stroke; all Q-family glyphs share a looped descender). There are
23 families with at least two members in the core inventory.

If the Voynich's positional structure is a designed code, a natural
prediction follows: \emph{glyphs that look alike should behave alike}.
That is, members of the same visual family should tend to occupy the
same word-internal positions. This is not trivially true: glyph
families are defined by graphic form, while word positions are defined
by distributional frequency. The correlation, if present, would
establish that the code's designer used visual similarity as a carrier
of functional information.

\subsection{Test design}

For each of the 23 families, we compute the Shannon entropy of its
members' positional distribution (the distribution of positions
$p_1$--$p_6$ occupied by all tokens in that family). The test statistic
is the mean family entropy across all 23 families. A lower mean entropy
indicates stronger positional concentration within families.

The null hypothesis is that visual family membership is unrelated to
positional behaviour. We implement a frequency-stratified permutation
test: glyph types are divided into frequency bands, and family labels
are permuted \emph{within} each band ($10{,}000$ iterations, random
seed 42). This controls for the trivial possibility that rare glyphs
are positionally restricted simply because they are rare.

\subsection{Results}

The observed mean family entropy is significantly lower than the
null distribution:

\begin{itemize}[nosep]
  \item Observed mean entropy: $H_{\text{obs}}$
  \item Null mean: $H_{\text{null}} = 2.213$~bits
  \item $z = -3.38$, $p = 0.0007$ (one-tailed)
\end{itemize}

The result is robust: the significance holds under both unstratified
and frequency-stratified null models, ruling out a frequency artefact.

We also tested a stronger hypothesis---that family membership predicts
\silpart\ positional \emph{class} (the six-class discretisation used
for generative modelling) rather than raw word position. This stronger
test is \emph{not} significant ($z = -0.11$, $p = 0.449$), indicating
that the correlation operates at the level of continuous positional
profiles, not discrete positional classes. Families predict
\emph{where in the word} a glyph tends to appear, but not the exact
\silpart\ slot boundary.

\subsection{Interpretation}

The family--position correlation establishes the foundational premise
for all subsequent analyses: \textbf{graphic form encodes functional
role in the Voynich code}. The designer of the system used visual
similarity systematically---not perfectly, but significantly above
chance---to signal which combinatorial slot a glyph belongs to. This
licenses the next step: classifying glyphs into functional roles based
on their positional profiles.


% ============================================================
\section{Six functional roles and combinatorial grammar}
\label{sec:role_grammar}

\subsection{Role assignment}

Building on the family--position correlation, we classify each of the
34 core glyphs (those with $\geq 100$ corpus occurrences, covering
99.5\% of all tokens) into functional roles based on their positional
concentration profile. The classification uses the dominant position
and the shape of the positional distribution to assign each glyph to
one of six roles:

\begin{enumerate}[nosep]
  \item \textbf{OPENER}: position-1 specialists. Highest frequency at
    $p_1$, rapid decay toward later positions. 10 glyph types, 37,533
    tokens. Examples: K1 (44.2\% at $p_1$), D1 (39.0\%), A1 (36.5\%).
  \item \textbf{ENCODER}: position-2 specialists. Highest frequency at
    $p_2$, encoding the primary category value. 7 glyph types, 57,220
    tokens.
  \item \textbf{CONNECTOR}: nuclear-dominant glyphs occupying
    mid-word positions ($p_3$--$p_4$). Internal structural function.
    3 glyph types, 17,381 tokens.
  \item \textbf{BODY}: late-position content carriers ($p_3$--$p_5$).
    3 glyph types, 21,608 tokens.
  \item \textbf{SCAFFOLD}: position-invariant structural glue,
    relatively flat distribution across positions. 3 glyph types,
    18,440 tokens.
  \item \textbf{CLOSER}: coda-dominant terminators, highest frequency
    at $p_5$--$p_6$. 8 glyph types, 3,136 tokens.
\end{enumerate}

\subsection{Transition matrix}

The six roles define a first-order Markov grammar over word-internal
transitions. We compute the empirical transition probability
$P(\text{role}_{i+1} \mid \text{role}_i)$ for all consecutive glyph
pairs within words across the entire corpus.

The grammar is strongly directional:

\begin{itemize}[nosep]
  \item OPENER $\to$ ENCODER: \textbf{63.5\%} of all position-1 to
    position-2 transitions.
  \item ENCODER $\to$ BODY: dominant at position-2 to position-3.
  \item BODY $\to$ CLOSER: dominant at late positions.
  \item Several role pairs are effectively forbidden ($O/E < 0.3$),
    including CLOSER $\to$ OPENER and SCAFFOLD $\to$ OPENER.
\end{itemize}

The transition matrix reveals a canonical word template:

\begin{center}
\textsc{OPENER} $\to$ \textsc{ENCODER} $\to$
[\textsc{BODY}/\textsc{CONNECTOR}] $\to$ \textsc{CLOSER}
\end{center}

\noindent with SCAFFOLD elements optionally inserted at any internal
position.

% TODO: Table 1 — Full 6×6 role transition matrix with O/E ratios
% (extract from notebook 76, cell 2)

\subsection{Section universality with local vocabulary}
\label{sec:sections}

The manuscript is traditionally divided into five thematic sections
based on illustration type: herbal, biological, astronomical, zodiac,
and pharmaceutical (recipe). If the combinatorial grammar is a property
of the encoding system rather than of any particular section's content,
we predict that the \emph{grammar} should be constant across sections
while the \emph{glyph proportions} vary.

This is precisely what we observe:

\begin{itemize}[nosep]
  \item The transition matrix is structurally stable across all five
    sections: OPENER $\to$ ENCODER remains the dominant first transition
    everywhere.
  \item What varies is the \emph{proportion} of specific glyphs within
    each role. Glyph D1 (an OPENER) is $1.7\times$ enriched in
    biological folios relative to the corpus baseline. Glyph Z1 is
    $2.8\times$ enriched in zodiac folios. These enrichments function
    as \emph{section markers}: glyphs whose elevated frequency in a
    particular section signals domain-specific content.
  \item A leave-one-out folio classifier using only the six role
    proportions per folio achieves \textbf{66.4\%} accuracy across the
    five sections, against a 27\% chance baseline ($p < 10^{-6}$).
\end{itemize}

The implication is that the code has a \emph{universal grammar} (the
transition rules) instantiated with \emph{local vocabularies} (the
section-specific glyph weightings). The grammar is the skeleton of the
code; the vocabulary is where the content lives.


% ============================================================
\section{Structural dictionary: the five-slot schema}
\label{sec:structural_dict}

\subsection{From roles to slots}

The six functional roles, combined with the positional structure
established by \silpart, induce a five-slot schema for the Voynich
word. Each slot corresponds to a positional range and a functional
interpretation derived from the role grammar:

\begin{enumerate}[nosep]
  \item \textbf{DOMINIUM} (position $p_1$): the domain classifier.
    Selects the broad thematic domain of the entry. Closed inventory
    of 113 glyph types (36,568 tokens). Highest section variability
    (Jensen--Shannon divergence $\text{JSD} = 0.226$). Top glyphs:
    A1 (22.1\%), K1 (16.0\%), D1 (14.6\%).
  \item \textbf{GENUS} (position $p_2$): the primary category encoder.
    Specifies the main category within the domain. 89 types (35,697
    tokens), $\text{JSD} = 0.189$. Top glyphs: A1 (26.2\%), J1
    (12.3\%), A3 (10.8\%), Q1 (10.2\%), Q2 (8.5\%).
  \item \textbf{SPECIES} (position $p_3$): the secondary modifier.
    Carries the highest informational load of any slot (Shannon entropy
    $H = 4.09$~bits). 75 types (32,232 tokens), $\text{JSD} = 0.199$.
    Top glyphs: Q1 (13.2\%), A3 (10.9\%), A1 (9.8\%), J1 (8.4\%).
  \item \textbf{CORPUS} (position $p_4$): body content. 81 types
    (24,152 tokens), $\text{JSD} = 0.186$ (lowest section variability,
    suggesting section-invariant structural function). Top glyphs:
    A2 (17.9\%), A3 (12.8\%), A1 (9.7\%), B2 (8.3\%).
  \item \textbf{FINIS} (positions $p_5$--$p_6$): the terminator.
    60 types at $p_5$ (14,962 tokens), 51 types at $p_6$ (6,804
    tokens). Glyph A2 dominates $p_6$ at 38.1\% with the lowest
    coefficient of variation ($\text{CV} = 0.09$), making it the most
    universal closer in the system---effectively a full stop.
    $\text{JSD} = 0.213$ at $p_6$ (second highest, suggesting that
    termination patterns carry section-specific information).
\end{enumerate}

\subsection{Slot properties}

Several properties of the schema deserve explicit comment.

\paragraph{Entropy gradient.}
The informational load is not uniformly distributed. Position $p_3$
(SPECIES) carries the highest entropy, followed by $p_2$ (GENUS) and
$p_1$ (DOMINIUM). Positions $p_5$--$p_6$ (FINIS) carry the lowest.
This gradient is consistent with a classificatory system: the first
slots narrow the domain progressively, and the terminal slots close
the entry with a constrained set of markers.

\paragraph{Universal closer.}
Glyph A2 at position $p_6$ is effectively a punctuation mark: it
appears in 38.1\% of all sixth-position tokens with negligible
section variation ($\text{CV} = 0.09$). No other glyph in any other
position achieves comparable universality.

\paragraph{Section specificity.}
The Jensen--Shannon divergence across the five manuscript sections
varies by slot. DOMINIUM ($p_1$) shows the highest JSD (0.226),
confirming that the opening glyph is where section-specific domain
selection occurs. CORPUS ($p_4$) shows the lowest (0.186), suggesting
that the body of the word is structurally stable regardless of section.

\subsection{Machine-readable export}

The full structural dictionary---including glyph inventories, frequency
profiles, section specificities, and slot assignments for all 166 glyph
types across six positions---is exported as a JSON file
(\filename{voynich_structural_dictionary.json}) available in the
project's public repository. This file provides the complete
specification of the reconstructed schema in machine-readable form,
suitable for programmatic validation or extension by other researchers.


% ============================================================
\section{CLAVIS CODICIS: the schema of the lost key book}
\label{sec:clavis}

The five-slot schema reconstructed in the previous section represents
the \emph{syntactic skeleton} of the document that is missing: the
original key book (or equivalent reference table) without which the
manuscript's content cannot be read.

We can now state precisely what the key book must contain and what
form it must take:

\begin{enumerate}[nosep]
  \item \textbf{Five assignment tables}, one per slot (DOMINIUM through
    FINIS), mapping each glyph in that slot's inventory to a semantic
    value in the original author's conceptual framework.
  \item \textbf{Section prefixes or markers} indicating how the
    assignment tables vary across the manuscript's five thematic
    sections---or, equivalently, five parallel sets of tables, one
    per section.
  \item \textbf{A reading rule} specifying the order in which slots
    are composed to produce a complete entry: DOMINIUM selects the
    domain, GENUS the category, SPECIES the modifier, CORPUS the
    body content, and FINIS terminates the entry.
\end{enumerate}

What the structural analysis \emph{cannot} recover is the semantic
content of the assignment tables: the actual meanings that the original
author assigned to each glyph in each slot. The grammar tells us that
glyph K1 at position~1 is a domain classifier enriched in herbal
folios; it does not tell us which herbal domain K1 designates. The
grammar tells us that the SPECIES slot carries the highest
informational load; it does not tell us what the species categories
are.

This is the fundamental asymmetry of the reconstruction: the analysis
recovers the \emph{rules} of the code but not its \emph{content}. The
syntax is accessible from the text alone; the semantics require the
key book. This asymmetry is not a limitation of the method---it is a
property of the encoding system itself. A categorial code is designed
so that its combinatorial structure can be learned (by the intended
reader, who has the key book) while its content remains opaque to
anyone who does not.

The reconstructed schema nonetheless has immediate practical value:
any future candidate for the lost key book must be \emph{compatible}
with the five-slot structure. A historical document that proposes
semantic assignments violating the grammar (e.g., assigning an OPENER
glyph to a terminator function, or proposing more category values
than the slot's inventory allows) can be ruled out on structural
grounds alone. The schema thus functions as a \emph{filter} for
documentary research: it does not tell us where the key book is, but
it tells us what shape it must have.

% TODO: Figure — CLAVIS CODICIS visualisation from notebook 78
% (key_book_reconstructed.png)


% ============================================================
\section{Discussion}
\label{sec:discussion}

The analysis presented in this paper recovers the combinatorial
structure of the Voynich code at a level of detail that, to our
knowledge, has not been previously achieved. The progression from
form--function correlation (\S\ref{sec:family_position}) through
role grammar (\S\ref{sec:role_grammar}) to structural dictionary
(\S\ref{sec:structural_dict}) yields a complete syntactic
specification of the encoding system: we know what the valid word
forms are, what roles each glyph plays, how roles combine, and how
the system varies across manuscript sections.

What the analysis does \emph{not} yield is the semantic content.
The structural dictionary delimits the space of valid Voynich words
but does not tell us what any of them mean. This is not a failure of
the method but a direct consequence of the encoding type: in a
categorial code, the syntax is public (it is visible in the statistical
regularities of the text) while the semantics are private (they are
stored in the key book, which is not part of the text). The analysis
can describe the rules of the system from outside without being able
to read the content from inside.

This asymmetry has a constructive implication. If the code's syntax
is now known, then the search for the lost key book is no longer
unconstrained. Any candidate document must satisfy the five-slot
schema: it must provide assignment tables for closed glyph inventories
at each position, it must respect the transition grammar, and it must
account for section-specific variation. This transforms the documentary
search from ``find a book that might explain the Voynich'' to ``find
a book whose structure matches a known five-slot categorial schema
with these specific properties.'' The search space is smaller than it
was before.

\subsection{Relation to prior work}

The positional structure of Voynich words has been noted by several
previous authors, most notably in the ``slot grammar'' descriptions
of the voynichological community and in the formal analyses of
\citet{dimperio1978}. The present work differs from prior descriptions
in three respects: (i)~it provides a statistical test for the
family--position correlation ($p = 0.0007$) rather than relying on
visual inspection; (ii)~it derives a complete six-role grammar with
quantified transition probabilities rather than a qualitative
positional description; and (iii)~it exports the result as a
machine-readable structural dictionary suitable for programmatic
use.

\subsection{Relation to the Picatrix framework}

The third companion paper~\citep{jimenez2026venusring} established
that the \vms\ belongs to the medieval astrological-talismanic
tradition codified by \textit{Picatrix Latinus}, and that the
manuscript's encoding constitutes a \textit{cifra ineffabilis} in the
sense recovered by \citet{martinezgavilan2025}. The present paper's
findings are consistent with that framework but do not depend on it:
the role grammar and structural dictionary are derived entirely from
the text's internal statistics and hold regardless of the manuscript's
historical attribution. If the Picatrix connection were refuted, the
five-slot schema would remain intact.


% ============================================================
\section{Limitations and future work}
\label{sec:limitations}

\paragraph{Role assignment granularity.}
The six-role classification is based on dominant positional profiles
of the 34 core glyphs. The remaining 132 glyph types (0.5\% of
tokens) are not classified. A finer-grained classification---for
instance, distinguishing sub-roles within ENCODER---might reveal
additional structure but would require a larger effective sample size
per glyph type than the corpus currently provides.

\paragraph{Transition order.}
The grammar reported here is first-order Markov (bigram transitions
between consecutive roles). Higher-order dependencies may exist but
are difficult to estimate reliably given the short average word
length (4.24 glyphs/word). Future work could explore trigram
transitions for the most frequent role sequences.

\paragraph{Section boundaries.}
The five-section partition (herbal, biological, astronomical, zodiac,
pharmaceutical) follows the traditional iconographic classification
and may not correspond exactly to the original author's thematic
divisions. The section-specificity analysis is sensitive to this
partition; alternative segmentations should be tested.

\paragraph{Slot nomenclature and interpretive framework.}
The Latin slot names (DOMINIUM, GENUS, SPECIES, CORPUS, FINIS) are
an \emph{interpretive} overlay, not an empirical result. What the
data establish is that five positional slots exist with closed
inventories, measurable entropy gradients, and section-specific
variation. The decision to label these slots with medieval scholastic
categories reflects the Picatrix/Caramuel framework adopted in the
companion papers; a critic working from a purely linguistic
perspective could relabel them as prefix, root, infix, suffix, and
terminator without altering any quantitative finding. The nomenclature
is chosen for historical coherence with the manuscript's probable
tradition, not because the statistics mandate a taxonomic rather than
grammatical reading. Readers should evaluate the structural results
independently of the labels assigned to them.

\paragraph{Semantic access.}
The structural dictionary provides the syntactic frame but not
the semantic content of the code. Content access requires either
the recovery of the lost key book or the identification of a
sufficiently constrained set of ``cribs'' (known plaintext--ciphertext
pairs) that could bootstrap semantic assignment. The zodiac labels
analysed in the third companion paper~\citep{jimenez2026venusring}
represent the most promising crib candidates identified to date, but
even they proved insufficient for cross-sign discrimination.

\paragraph{Reproducibility.}
All analyses are implemented in Jupyter notebooks (notebooks 75--78
in the project repository), using the public \sta\ 2.0 corpus. The
structural dictionary is available as a JSON file. All code and data
required to reproduce every figure and table in this paper are
publicly available under the project's Zenodo DOI.


% ============================================================
\section{Conclusion}
\label{sec:conclusion}

We have shown that the combinatorial rules of the \vms's encoding
system can be substantially recovered from the text alone. Visual
glyph families correlate significantly with word-internal positions
($z = -3.38$, $p = 0.0007$), resolving into a six-role grammar with
directional transition rules (OPENER $\to$ ENCODER at 63.5\%) that
is universal across all five manuscript sections. The grammar induces
a five-slot structural schema (DOMINIUM, GENUS, SPECIES, CORPUS,
FINIS) with closed inventories and measurable section specificity,
representing the syntactic skeleton of the lost key book.

The reconstruction establishes that the Voynich code is not
opaque because it lacks structure---it is opaque because its structure
is \emph{syntactic}, and the missing piece is \emph{semantic}. The
rules are in the text. The meanings are in the key book. The
structural dictionary now provides a formal specification against
which any future key-book candidate can be tested.

The manuscript's content remains unread. But the shape of the lock
is now known, and that constrains the shape of the key.
